\section{The intermediate range}\label{sec:discussion}
The answer is positive for every $N\le50$ and negative for every
$N\ge\ObstructionOrder$. By Theorem~\ref{thm:exchange}, the remaining
orders are decided by the sign of $\gamma_N-1$, which the present
estimates do not determine. Theorem~\ref{thm:criticality} forces the
relative pure margins of any certificates toward zero as the order
grows, and Theorem~\ref{thm:negative} gives $\gamma_N>1$ for
$N\ge\ObstructionOrder$; neither result locates a transition.

\begin{remark}[An expectation at moderate orders]\label{rem:persistence}
The exact positive certificates through order $50$ and the small
normalized defect $2^{-18}<\delta_{\rm test}<2^{-16}$ of the tangent
witness (Appendix~\ref{app:witness}) suggest that the positive SOS
range may extend to orders in the hundreds. This is an expectation
without a specified endpoint, not a proved extension of the positive
range. The continuation estimate explains why the present obstruction
becomes effective only much later; it cannot exclude a different
obstruction at a smaller order.
\end{remark}

\subsection{Open questions}\label{sec:questions}
Is there an integer $N_*$ such that $F_N$ is a sum of squares exactly
for the orders $1\le N<N_*$? If so, $51\le N_*\le\ObstructionOrder$.
A first failure exists by Theorem~\ref{thm:negative}; the question is
whether the sum-of-squares property can return at a larger order. No
monotonicity theorem is known for $\gamma_N$ or for the forms $F_N$.

Further questions concern the first non-SOS order, improvements of the
sufficient threshold through a different tangent test or a stronger
continuation estimate, and exact positive certificates at intermediate
orders.
